\documentclass[11pt]{article}
\usepackage[a4paper,margin=2.6cm]{geometry}
\usepackage{amsmath,amssymb,amsthm}
\usepackage{hyperref}
\newtheorem{prop}{Proposition}
\newcommand{\dd}{\mathrm{d}}
\title{The sign of $\sqrt{xyz}$ as a Galois ($\mathbb{Z}_2$) symmetry\\ of the canonical differential equation}
\author{Notes for the H family (3L1M), analytic continuation to $x,y<0$}
\date{3 October 2026}
\begin{document}
\maketitle

\section{Setup}
The masters $\vec g(x,y;\epsilon)$ of the H family ($N_f=371$) satisfy a canonical
differential equation
\begin{equation}
  \dd \vec g = \epsilon\, \dd A\, \vec g ,\qquad
  A = \sum_{k} a_k \log l_k(x,y), \qquad a_k\in\mathbb{Q}^{N_f\times N_f},
\end{equation}
with $\vec g = \sum_{w\ge0} \epsilon^w \vec g^{(w)}$. The alphabet contains one square root,
\begin{equation}
  r=\sqrt{b},\qquad b = x\,y\,z,\qquad z = 1-x-y ,
\end{equation}
and lives in the quadratic field extension $K(r)$ of $K=\mathbb{Q}(x,y)$.
The letters split into two classes:
\begin{itemize}
\item \emph{even} (rational) letters $l_1,\dots,l_6,l_9,\dots,l_{20}\in K$;
\item \emph{odd} letters
\begin{equation}
  l_7=\frac{x z - i r}{x z + i r},\qquad l_8=\frac{x y - i r}{x y + i r}.
\end{equation}
\end{itemize}

\section{The Galois action}
$\mathrm{Gal}(K(r)/K)=\{1,\sigma\}$ with $\sigma: r\mapsto -r$. On the letters,
\begin{equation}
  \sigma(l_k)=l_k \ (\text{even}),\qquad \sigma(l_{7,8}) = l_{7,8}^{-1}
  \;\Rightarrow\; \sigma(\dd\log l_{7,8}) = -\dd\log l_{7,8} .
\end{equation}
Choosing the sign of $r$ at a point (e.g.\ after an analytic continuation) is choosing
between $1$ and $\sigma$.

The masters split in the same way. In H the only masters normalized with a square-root
prefactor are $g_{156}, g_{157}$ (``odd''); all others are ``even''. Write
$P=\mathrm{diag}(1,\dots,1,-1,-1,1,\dots,1)$ with $-1$ on the odd masters.

\section{Parity structure of the connection}
Write $A$ in even/odd blocks,
\begin{equation}
  A=\begin{pmatrix} A_{ee} & A_{eo}\\ A_{oe} & A_{oo}\end{pmatrix}.
\end{equation}
The DE is \emph{graded}: the diagonal blocks contain only even letters, and the off-diagonal
blocks contain only odd letters. For H this was checked entry by entry on
\texttt{H\_atilde}: every one of the 134 nonzero odd$\leftrightarrow$even entries contains only
$l_7,l_8$, and $l_7,l_8$ appear nowhere else. Equivalently,
\begin{equation}
  \sigma(A) = P\,A\,P .
  \label{eq:cov}
\end{equation}
Note that the odd--odd entries (e.g.\ $A_{156,156}$) do \emph{not} change sign: they multiply an
odd master and contribute to the derivative of an odd master, so the two signs cancel.

\section{Consequence}
\begin{prop}
Let $\vec g$ solve $\dd\vec g=\epsilon\,\dd A\,\vec g$ with boundary value $\vec g_0$. Then
$P\vec g$ solves the Galois-conjugate equation $\dd\vec g'=\epsilon\,\dd\sigma(A)\,\vec g'$
with boundary value $P\vec g_0$. If the odd boundary constants vanish, $P\vec g_0=\vec g_0$,
so the solution of the conjugate equation with the same boundary data is $P\vec g$:
\begin{equation}
  g_e\ \text{unchanged},\qquad g_{156},\,g_{157}\ \to\ -g_{156},\,-g_{157}.
\end{equation}
\end{prop}
\begin{proof}
$\dd(P\vec g)=\epsilon\,P\,\dd A\,\vec g=\epsilon\,(P\,\dd A\,P)(P\vec g)=\epsilon\,\dd\sigma(A)\,(P\vec g)$,
using $P^2=1$ and \eqref{eq:cov}.
\end{proof}
Block by block:
\begin{align}
  \dd g_e &= \epsilon\,(\dd A_{ee}\, g_e + \dd A_{eo}\, g_o), &
  \dd(-g_o) &= \epsilon\,(\dd A_{oo}\,(-g_o) + (-\dd A_{oe})\, g_e),
\end{align}
i.e.\ in the even rows the flipped odd letter multiplies the flipped odd master, and the odd
rows are the original equation times $-1$.

\paragraph{Symbol-level statement.} Since $g^{(w)}=\int \dd A\, g^{(w-1)} + c^{(w)}$ and the odd
constants $c^{(w)}_{o}$ vanish (including weight 0), induction on the weight shows that every
term in the symbol (iterated integral) of an even master contains an even number of odd
letters, and every term of an odd master an odd number. Under $\sigma$, odd masters are odd and even masters are invariant;
the masters are Galois eigenvectors with eigenvalue $\pm1$.

\section{Application to the continued region}
In the chart $x=-w/v$, $y=-u/v$, $z=1/v$ ($u+w+v=1$), which covers $x,y<0$, $z>1$, one has
$xyz = uw/v^3>0$, so $r$ is real, but its sign is a choice. For H:
\begin{itemize}
\item the boundary constants of $g_{156},g_{157}$ vanish at every weight at the old corners
  c1, c2, c3 and at the continued points CC1, CC2, CC3;
\item hence the choice $r=+\sqrt{|xyz|}$ versus $r=-\sqrt{|xyz|}$ changes only the overall sign
  of $g_{156},g_{157}$; the other 369 masters are identical;
\item the choice consistent with $\log x\to\log|x|+i\pi$, $\log y\to\log|y|+i\pi$ (the
  continuation used for the CC boundaries) is
  $r=\sqrt{x}\sqrt{y}\sqrt{z}=e^{i\pi}\sqrt{|xyz|}=-\sqrt{|xyz|}$.
  It matters only when $g_{156},g_{157}$ are combined with quantities defined in the original
  region (e.g.\ the prefactor $\sqrt{xyz}$ itself).
\end{itemize}
The letter expansions at CC2, CC3 (\texttt{cc\_series/letters/build\_cc\_letters.py}) use
$r=-\sqrt{|xyz|}$ (default \texttt{--odd-sign -1}), consistent with the $+i\pi$ continuation of
$\log x$ and $\log y$; \texttt{--odd-sign +1} gives the principal root.

\paragraph{Where the sign comes from.} The radicand is positive at both ends of the
continuation ($xyz>0$ for $x,y>0$ and for $x,y<0$, $z>1$), but negative in between (e.g.\ $x<0<y$),
where $r$ is imaginary. Each crossing contributes $e^{\pm i\pi/2}$:
$\sqrt{x}\to i\sqrt{|x|}$, $\sqrt{y}\to i\sqrt{|y|}$ for $\log x,\log y\to\log|\cdot|+i\pi$, so
$r\to -\sqrt{|xyz|}$: a minus sign in front of the root, not under it. Opposite phases for $x$ and
$y$ would give $+\sqrt{|xyz|}$.

\end{document}
